\begin{claim}{D4.2}{Proved}{The cusp length and the generic twist}
Use the coordinate $M_z=\exp(z\mathfrak g)$ of C3.2. For every
$z\notin\Z$ the cusp contribution is exactly
\[
 (R^0f_*(T_X\otimes M_z))_{p_0}=0,\quad
 (R^1f_*(T_X\otimes M_z))_{p_0}=0,\quad
 (R^2f_*(T_X\otimes M_z))_{p_0}=\C\{q\}/(q).           \tag{D.14}
\]
Outside the countable exceptional set explicitly defined in D4.3,
these are the only nonzero direct-image stalks. Hence the generic
twist has $(h^0,h^1,h^2,h^3)=(0,0,1,0)$, and the integer $r$ in
Lead 5 equals one there.
\end{claim}
\begin{proof}
The map $H^1(W,\C)\to H^1(W,\OO_W)$ is an isomorphism by the
normalization complex in D1.2. The restriction of $M_z$ has trivial
pullback to $P$ and descent character
$(\exp(2\pi iz),1)$ on the two generators of the dual torus. The
normalization complex with these coefficients is the cellular
cochain complex of that rank-one local system. It may be computed
on the product cell structure as
\[
 \C\xrightarrow{(a-1,0)^t}\C^2
 \xrightarrow{(0,a-1)}\C,
 \qquad a=\exp(2\pi iz).
\]
For $z\notin\Z$ this is acyclic. All normalization components remain
rational with trivial pulled-back twist, so this calculates the
actual coherent cohomology $H^q(W,M_z|_W)$, not just a topological
analogy. Analytic semicontinuity and base change make
$R^q\bar f_*M_z$ zero on a neighbourhood of $p_0$ for all $q$.
The projection formula and (D.3) imply the same vanishing for
$\mathcal T\otimes M_z$. Tensoring (D.10) remains exact, and
$\nu^*M_z=\OO_P$ because $P$ is rational and the topological first
Chern class of $M_z$ is zero. Thus the degree-two direct image is
exactly the reduced skyscraper $H^2(P,K_P)\otimes\C_{p_0}$.
This proves (D.14), including its annihilator and not only its residue
dimension.

If the twist is nontrivial on every smooth fibre and on both finite
root-cover central tori, all cohomology of its fibre restriction
vanishes at those places. Indeed a nontrivial topologically trivial
line bundle on a complex torus has zero cohomology: in a unitary flat
presentation its Dolbeault Laplacian has no zero Fourier mode.
Tensoring the tangent extension by that line preserves acyclicity.
Base change on the root charts and exact invariants then give zero
direct images away from $p_0$. D4.3 proves that the excluded twists
form a countable set. Leray now gives the asserted four dimensions.
\end{proof}
